\section{VAE: variational autoencoder}

\subsection{Motivación de pasar de GAN a VAE}

La dificultad central de los GAN radica en la optimización minimax. La idea detrás del VAE es: \textbf{¿podemos evitar el entrenamiento adversarial y aprender la mapeo de la distribución latente a la distribución de datos mediante una forma más estable?}

La innovación clave del VAE consiste en dejar de aprender una función de mapeo determinista $G_\theta(z)$ para aprender una \textbf{distribución de probabilidad condicional} $p_\theta(x|z)$, es decir, la distribución de los datos $x$ dado que tenemos la variable latente $z$.

\subsection{Marco probabilístico del VAE}

\subsubsection{Proceso generativo}

El VAE define el siguiente proceso generativo:

\begin{enumerate}
    \item Muestreo desde una distribución a priori: $z \sim p(z) = \mathcal{N}(\mathbf{0}, \mathbf{I})$
    \item Muestreo desde la distribución de verosimilitud: $x \sim p_\theta(x|z) = \mathcal{N}(\mu_\theta(z), \sigma^2 \mathbf{I})$
\end{enumerate}

Donde $\mu_\theta(z)$ es la salida de la red neuronal del decodificador (la predicción del valor medio de $x$ dado $z$) y $\sigma^2$ es un hiperparámetro de varianza fija.

\begin{importantbox}{Modelo probabilístico generativo del VAE}
$$p_\theta(x, z) = p_\theta(x|z) \cdot p(z)$$
Donde:
\begin{itemize}
    \item $p(z) = \mathcal{N}(\mathbf{0}, \mathbf{I})$: distribución a priori (prior), conocida.
    \item $p_\theta(x|z) = \mathcal{N}(\mu_\theta(z), \sigma^2 \mathbf{I})$: distribución de verosimilitud (likelihood), parametrizada por el decodificador.
    \item $p_\theta(x) = \int p_\theta(x|z) p(z) \, dz$: verosimilitud marginal (marginal likelihood), obtenida integrando respecto a $z$.
\end{itemize}
\end{importantbox}

\subsubsection{Objetivo de entrenamiento: maximizar la verosimilitud marginal}

En principio, deseamos maximizar el logaritmo de la verosimilitud marginal de los datos:

$$\max_\theta \; \mathbb{E}_{x \sim p_{\text{data}}} [\log p_\theta(x)] = \max_\theta \; \mathbb{E}_{x \sim p_{\text{data}}} \left[\log \int p_\theta(x|z) p(z) \, dz \right]$$

\begin{warningbox}{La verosimilitud marginal no es calculable}
La integral $p_\theta(x) = \int p_\theta(x|z) p(z) \, dz$ \textbf{no puede calcularse de forma exacta} en espacios de alta dimensión. Intuitivamente: la dimensión del espacio latente podría ser de cientos de dimensiones, e integrar sobre todos los posibles $z$ no es factible. Por eso necesitamos inferencia variacional (Variational Inference).
\end{warningbox}

\subsection{Repaso de probabilidad: herramientas matemáticas necesarias para el VAE}

\subsubsection{Marginalización}

Dada una distribución conjunta $p(x, z)$, se puede obtener la distribución marginal de otra variable integrando (o sumando) respecto a una de ellas:

$$p(x) = \int p(x, z) \, dz = \int p(x|z) p(z) \, dz$$

\subsubsection{Teorema de Bayes}

$$p(z|x) = \frac{p(x|z) p(z)}{p(x)}$$

Convención de terminología para cada término:

\begin{table}[H]
\centering
\begin{tabular}{ll}
\toprule
Símbolo & Nombre \\
\midrule
$p(z)$ & Prior (a priori): conocimiento a priori sobre $z$ \\
$p(x|z)$ & Likelihood (verosimilitud): probabilidad de observar $x$ dado $z$ \\
$p(z|x)$ & Posterior (a posteriori): creencia actualizada sobre $z$ tras observar $x$ \\
$p(x)$ & Evidence / Marginal Likelihood (evidencia): probabilidad marginal de los datos \\
\bottomrule
\end{tabular}
\caption{Terminología de los términos en el teorema de Bayes}
\end{table}

\subsubsection{Divergencia de Kullback-Leibler}

La divergencia de Kullback-Leibler mide la "distancia" entre dos distribuciones de probabilidad (nota: no es simétrica, por lo que no es estrictamente una métrica de distancia):

$$D_{\mathrm{KL}}(p \| q) = \mathbb{E}_{x \sim p} \left[\log \frac{p(x)}{q(x)}\right] = \int p(x) \log \frac{p(x)}{q(x)} \, dx$$

\begin{importantbox}{Propiedades clave de la divergencia de KL}
\begin{itemize}
    \item \textbf{No negatividad}: $D_{\mathrm{KL}}(p \| q) \geq 0$, con igualdad si y solo si $p = q$ casi en todas partes.
    \item \textbf{Asimetría}: generalmente $D_{\mathrm{KL}}(p \| q) \neq D_{\mathrm{KL}}(q \| p)$.
    \item \textbf{Forma cerrada para la divergencia KL entre dos distribuciones gaussianas}.
\end{itemize}
\end{importantbox}

\begin{knowledgebox}{Divergencia KL entre dos distribuciones gaussianas multivariadas}
Supongamos $p = \mathcal{N}(\boldsymbol{\mu}_1, \boldsymbol{\Sigma}_1)$ y $q = \mathcal{N}(\boldsymbol{\mu}_2, \boldsymbol{\Sigma}_2)$. Entonces:
$$D_{\mathrm{KL}}(p \| q) = \frac{1}{2}\left[\log\frac{|\boldsymbol{\Sigma}_2|}{|\boldsymbol{\Sigma}_1|} - d + \mathrm{tr}(\boldsymbol{\Sigma}_2^{-1}\boldsymbol{\Sigma}_1) + (\boldsymbol{\mu}_2 - \boldsymbol{\mu}_1)^T \boldsymbol{\Sigma}_2^{-1} (\boldsymbol{\mu}_2 - \boldsymbol{\mu}_1)\right]$$
donde $d$ es la dimensión. En particular, cuando $q = \mathcal{N}(\mathbf{0}, \mathbf{I})$:
$$D_{\mathrm{KL}}(\mathcal{N}(\boldsymbol{\mu}, \boldsymbol{\Sigma}) \| \mathcal{N}(\mathbf{0}, \mathbf{I})) = \frac{1}{2}\left[-\log|\boldsymbol{\Sigma}| - d + \mathrm{tr}(\boldsymbol{\Sigma}) + \boldsymbol{\mu}^T\boldsymbol{\mu}\right]$$
Esta fórmula se usa directamente en el entrenamiento del VAE.
\end{knowledgebox}

\subsection{Derivación de la desigualdad de Jensen y el ELBO}

\subsubsection{Funciones convexas e inecuación de Jensen}

\textbf{Función convexa (Convex Function)}: se define como aquella que, para cualquier par de puntos $x_1, \ldots, x_n$ y pesos $t_1, \ldots, t_n$ ($t_i \geq 0$, $\sum t_i = 1$), cumple:

$$f\left(\sum_i t_i x_i\right) \leq \sum_i t_i f(x_i)$$

Intuitivamente, el gráfico de una función convexa se curva "hacia abajo", y cualquier segmento que une dos puntos yace sobre el gráfico de la función.

Una \textbf{función cóncava (Concave Function)} es lo opuesto a la convexa: su gráfico se curva "hacia arriba" y $-f$ es una función convexa. La función $\log$ es una función cóncava importante.

\begin{importantbox}{Inecuación de Jensen}
Si $f$ es una función convexa y $X$ es una variable aleatoria, entonces:
$$f(\mathbb{E}{[}X{]}) \leq \mathbb{E}{[}f(X){]}$$
Si $f$ es cóncava (como $\log$), la dirección de la inecuación se invierte:
$$f(\mathbb{E}{[}X{]}) \geq \mathbb{E}{[}f(X){]}$$
Es decir, $\log(\mathbb{E}{[}X{]}) \geq \mathbb{E}{[}\log X{]}$.
\end{importantbox}

\subsubsection{Derivación del ELBO}

El ELBO (Evidence Lower Bound) es el objetivo central del entrenamiento del VAE. A continuación se muestra el proceso de derivación completo.

\textbf{Punto de partida}: queremos maximizar el logaritmo de la verosimilitud marginal $\log p_\theta(x)$, pero no puede calcularse directamente. Introduciendo una distribución a posteriori aproximada $q_\phi(z|x)$ (encoder), tenemos:

\begin{align}
\log p_\theta(x) &= \log \int p_\theta(x, z) \, dz \nonumber \\
&= \log \int \frac{p_\theta(x, z)}{q_\phi(z|x)} q_\phi(z|x) \, dz \nonumber \\
&= \log \mathbb{E}_{z \sim q_\phi(z|x)} \left[\frac{p_\theta(x, z)}{q_\phi(z|x)}\right] \nonumber \\
&\geq \mathbb{E}_{z \sim q_\phi(z|x)} \left[\log \frac{p_\theta(x, z)}{q_\phi(z|x)}\right] \quad \text{(Inecuación de Jensen, $\log$ es cóncava)}
\end{align}

La última línea es el ELBO:

$$\text{ELBO}(\theta, \phi; x) = \mathbb{E}_{z \sim q_\phi(z|x)} \left[\log \frac{p_\theta(x, z)}{q_\phi(z|x)}\right]$$

\begin{importantbox}{Descomposición del ELBO}
Sustituyendo $p_\theta(x, z) = p_\theta(x|z)p(z)$, el ELBO se puede descomponer en dos términos:
\begin{align}
\text{ELBO} &= \mathbb{E}_{z \sim q_\phi(z|x)} \left[\log p_\theta(x|z)\right] - D_{\mathrm{KL}}(q_\phi(z|x) \| p(z)) \nonumber
\end{align}
\begin{itemize}
    \item \textbf{Término de reconstrucción (Reconstruction Term)}: $\mathbb{E}_{z \sim q_\phi(z|x)} {[}\log p_\theta(x|z){]}$ —mide la capacidad del decodificador para reconstruir $x$ a partir de $z$. Cuando $p_\theta(x|z) = \mathcal{N}(\mu_\theta(z), \sigma^2 I)$, este término es equivalente al error cuadrático medio (MSE) negativo.
    \item \textbf{Término de regularización (Regularization / KL Term)}: $-D_{\mathrm{KL}}(q_\phi(z|x) \| p(z))$ —incentiva que la distribución de salida del encoder se acerque a la a priori $p(z) = \mathcal{N}(\mathbf{0}, \mathbf{I})$, estructurando así el espacio latente.
\end{itemize}
\end{importantbox}

\subsubsection{Relación entre ELBO y verosimilitud marginal}

En realidad, se puede escribir con más precisión:

$$\log p_\theta(x) = \text{ELBO}(\theta, \phi; x) + D_{\mathrm{KL}}(q_\phi(z|x) \| p_\theta(z|x))$$

\begin{warningbox}{El ELBO es una cota inferior, no una igualdad}
Dado que $D_{\mathrm{KL}}(q_\phi(z|x) \| p_\theta(z|x)) \geq 0$, el ELBO siempre será menor o igual al logaritmo de la verosimilitud marginal real $\log p_\theta(x)$. La diferencia (gap) es igual a la divergencia KL entre la a posteriori aproximada $q_\phi(z|x)$ y la a posteriori real $p_\theta(z|x)$. Cuando $q_\phi(z|x) = p_\theta(z|x)$ (la a posteriori aproximada coincide exactamente con la real), el ELBO es igual a $\log p_\theta(x)$ y el gap es cero.
\end{warningbox}

\subsection{Arquitectura de la red del VAE}

En términos de implementación, el VAE incluye:

\begin{itemize}
    \item \textbf{Red Encoder} $q_\phi(z|x)$: entrada $x$, salida $\mu_\phi(x)$ y $\sigma_\phi(x)$ (media y desviación estándar de la distribución latente).
    \item \textbf{Red Decoder} $p_\theta(x|z)$: entrada $z$, salida $\mu_\theta(z)$ (valor medio para reconstruir $x$).
    \item \textbf{Proceso de muestreo}: se realiza mediante el truco de reparametrización para obtener $z = \mu_\phi(x) + \sigma_\phi(x) \odot \epsilon$, con $\epsilon \sim \mathcal{N}(\mathbf{0}, \mathbf{I})$.
\end{itemize}

\begin{knowledgebox}{Papel clave del truco de reparametrización en el VAE}
La operación de muestrear $z$ directamente desde $q_\phi(z|x) = \mathcal{N}(\mu_\phi(x), \sigma^2_\phi(x) \mathbf{I})$ no es diferenciable y no permite optimizar $\phi$ mediante retropropagación. El truco de reparametrización expresa el muestreo como $z = \mu_\phi(x) + \sigma_\phi(x) \odot \epsilon$ (donde $\epsilon \sim \mathcal{N}(\mathbf{0}, \mathbf{I})$ es independiente de los parámetros), permitiendo que el gradiente fluya a través de $\mu_\phi$ y $\sigma_\phi$, logrando así un entrenamiento end-to-end.
\end{knowledgebox}

\subsection{Función de pérdida del VAE}

Sintetizando las derivaciones anteriores, el objetivo del entrenamiento del VAE es maximizar el ELBO, lo cual equivale a minimizar la siguiente función de pérdida:

$$\mathcal{L}_{\text{VAE}}(\theta, \phi; x) = -\text{ELBO} = -\mathbb{E}_{z \sim q_\phi(z|x)}{[}\log p_\theta(x|z){]} + D_{\mathrm{KL}}(q_\phi(z|x) \| p(z))$$

Cuando $p_\theta(x|z) = \mathcal{N}(\mu_\theta(z), \sigma^2 \mathbf{I})$ y $p(z) = \mathcal{N}(\mathbf{0}, \mathbf{I})$:

$$\mathcal{L}_{\text{VAE}} = \frac{1}{2\sigma^2}\|x - \mu_\theta(z)\|^2 + D_{\mathrm{KL}}(\mathcal{N}(\mu_\phi(x), \text{diag}(\sigma^2_\phi(x))) \| \mathcal{N}(\mathbf{0}, \mathbf{I}))$$

\begin{importantbox}{Intuición detrás de la función de pérdida del VAE}
La pérdida del VAE consta de dos partes:
\begin{enumerate}
    \item \textbf{Error de reconstrucción (MSE)}: obliga al modelo a reconstruir con precisión los datos de entrada.
    \item \textbf{Penalización KL}: fuerza a que la distribución latente se acerque a una gaussiana estándar, haciendo que el espacio latente sea suave, continuo y muestreable.
\end{enumerate}
Existe tensión entre ambos términos: el error de reconstrucción desea que las codificaciones latentes conserven al máximo la información (lo que podría llevar a distribuciones latentes superpuestas para diferentes puntos de datos), mientras que el término KL empuja a todas las distribuciones latentes hacia la gaussiana estándar (lo que podría provocar pérdida de información). Encontrar un buen equilibrio es el núcleo del ajuste de hiperparámetros en el VAE.
\end{importantbox}

\subsection{Relación entre el VAE y los modelos de difusión}

Durante la explicación del VAE, el profesor de clase adelantó varias veces: \textbf{la idea del VAE está directamente relacionada con los Modelos de Difusión}. En concreto:

\begin{itemize}
    \item Los Modelos de Difusión pueden entenderse como un tipo de VAE "jerárquico", donde las variables latentes son una serie de estados intermedios con ruido progresivo.
    \item El método de derivación del ELBO puede generalizarse directamente a los Modelos de Difusión.
    \item El truco de reparametrización presente en el VAE es igualmente indispensable en los Modelos de Difusión.
\end{itemize}

\subsection{Resumen de la sección}

El VAE evita las dificultades de optimización minimax de los GAN mediante inferencia variacional. Introduce un encoder (a posteriori aproximada) y un decodificador (verosimilitud), entrenando maximizando el ELBO (cota inferior del logaritmo de la verosimilitud marginal). El ELBO se descompone naturalmente en un término de reconstrucción y un término de regularización KL; el equilibrio entre ambos determina tanto la calidad generativa como el grado de estructuración del espacio latente. El entrenamiento del VAE es considerablemente más estable que el de los GAN, y su marco probabilístico sienta las bases teóricas para los posteriores Modelos de Difusión.

\newpage
